%! Special Directions --- Setting Up Eigenvalues — Unit 6, Lesson 6.0 — Guided Notes
\documentclass[10pt]{article}
\usepackage{linalg-article}
\usepackage{linalg-boxes}
\newcommand{\TallMath}[1]{$\displaystyle #1\rule[-1.4em]{0pt}{3.2em}$}
\newcommand{\xx}{\mathbf{x}}
\newcommand{\yy}{\mathbf{y}}
\newcommand{\vv}{\mathbf{v}}
\newcommand{\zero}{\mathbf{0}}

\begin{document}

\pageheader{Unit 6, Lesson 6.0}{Guided Notes}
\namedateperiod

\vspace{0.12in}

\begin{objectivebox}
\textbf{By the end of this lesson, I will be able to\dots}
\begin{itemize}\setlength{\itemsep}{2pt}
  \item compute $A\xx$ and decide whether $\xx$ is an \emph{eigenvector} (does $A\xx$ stay on the same line?);
  \item state $A\xx=\lambda\xx$ and read the \emph{eigenvalue} $\lambda$ as the stretch factor;
  \item explain why $\lambda=0$ means the matrix is \emph{not invertible} ($\det A=0$).
\end{itemize}
\end{objectivebox}

\vspace{0.08in}

\begin{vocabbox}
\small
Fill in each term as we build it together.
\vspace{2pt}
\termblanklong{Eigenvector (a special direction)}
\termblanklong{Eigenvalue $\lambda$ (the stretch factor)}
\termblanklong{The eigen-relation $A\xx=\lambda\xx$}
\termblanklong{$\lambda=0$ and ``not invertible'' (recall §5)}
\end{vocabbox}

\begin{hookbox}
A matrix multiplies each vector in the plane --- and usually \emph{turns} it to a new direction.
\begin{itemize}\setlength{\itemsep}{1pt}
  \item Are there directions a matrix does \emph{not} turn --- ones it only \emph{stretches}? \writeline
  \item If $A\xx$ points the same way as $\xx$ but is $3$ times as long, how would you write that with a symbol? \writeline
\end{itemize}
\end{hookbox}

\begin{notesbox}{1. A Matrix Usually Changes Direction}
Multiplying by a matrix moves a vector (Lesson~5.3). Take
$A=\begin{bsmallmatrix} 2 & 1\\ 1 & 2 \end{bsmallmatrix}$ and the vector
$\begin{bsmallmatrix} 1\\ 0 \end{bsmallmatrix}$:
\[
  A\begin{bmatrix} 1\\ 0 \end{bmatrix}
  = \begin{bmatrix} 2 & 1\\ 1 & 2 \end{bmatrix}\begin{bmatrix} 1\\ 0 \end{bmatrix}
  = \begin{bmatrix}\rule{0pt}{1.1em}\blank{0.9cm}\\[2pt]\blank{0.9cm}\end{bmatrix}.
\]
The output points in a \blank{2.4cm} direction than the input --- most vectors get turned. Today we
hunt for the exceptions.
\end{notesbox}

\begin{notesbox}{2. Special Directions: $A\xx=\lambda\xx$}
A nonzero vector $\xx$ is an \textbf{eigenvector} of $A$ when $A\xx$ lands on the \textbf{same line}
as $\xx$: it is only \emph{stretched}, not turned. In symbols,
\[
  A\xx = \lambda\xx \qquad (\xx\ne\zero),
\]
where the number $\lambda$ (``lambda'') is the \textbf{eigenvalue} --- the \emph{stretch factor}.
Test $\xx=\begin{bsmallmatrix} 1\\ 1 \end{bsmallmatrix}$ with the same $A$:
\[
  A\begin{bmatrix} 1\\ 1 \end{bmatrix}
  = \begin{bmatrix} 3\\ 3 \end{bmatrix}
  = \blank{0.9cm}\cdot\begin{bmatrix} 1\\ 1 \end{bmatrix},
  \quad\text{so } \lambda = \blank{0.9cm}.
\]

\begin{center}
\begin{tikzpicture}[scale=0.9]
  \draw[step=1, linegray!45, thin] (-0.4,-0.4) grid (3.5,3.5);
  \draw[->, thick, charcoal] (-0.5,0)--(3.7,0);
  \draw[->, thick, charcoal] (0,-0.5)--(0,3.7);
  % the eigen-line (same direction, only stretched)
  \draw[burgundy!45, dashed, thin] (-0.3,-0.3)--(3.4,3.4);
  \draw[->, ultra thick, burgundy] (0,0)--(1,1) node[below right=1pt]{\scriptsize $\xx=\begin{bsmallmatrix} 1\\1 \end{bsmallmatrix}$};
  \draw[->, very thick, burgundy!60] (0,0)--(3,3) node[above left=1pt]{\scriptsize $A\xx=3\xx$};
  % a non-eigenvector: turned
  \draw[->, ultra thick, royalblue] (0,0)--(1,0) node[below=2pt]{\scriptsize $\yy$};
  \draw[->, very thick, royalblue!55, dashed] (0,0)--(2,1) node[right=1pt]{\scriptsize $A\yy$ (turned)};
\end{tikzpicture}
\end{center}
Along the dashed line, $A$ only \emph{scales} $\xx$ (here $\times3$); the blue vector $\yy$ is
\blank{2.0cm} to a new direction. Now test $\xx=\begin{bsmallmatrix} 1\\ -1 \end{bsmallmatrix}$:
\[
  A\begin{bmatrix} 1\\ -1 \end{bmatrix}
  = \begin{bmatrix}\rule{0pt}{1.1em}\blank{0.9cm}\\[2pt]\blank{0.9cm}\end{bmatrix}
  = \blank{0.9cm}\cdot\begin{bmatrix} 1\\ -1 \end{bmatrix},
  \quad\text{so } \lambda = \blank{0.9cm}.
\]
\end{notesbox}

\begin{notesbox}{3. How to Test a Candidate}
Given $A$ and a vector $\xx$: \textbf{compute $A\xx$}, then ask ``\emph{is it a multiple of $\xx$?}''
If $A\xx=\lambda\xx$ for some number $\lambda$, then $\xx$ is an eigenvector and $\lambda$ is its
eigenvalue. If $A\xx$ is \emph{not} a multiple of $\xx$, it is not an eigenvector. Check
$\xx=\begin{bsmallmatrix} 2\\ 1 \end{bsmallmatrix}$:
\[
  A\begin{bmatrix} 2\\ 1 \end{bmatrix}
  = \begin{bmatrix} 5\\ 4 \end{bmatrix}.
\]
Is $\begin{bsmallmatrix} 5\\ 4 \end{bsmallmatrix}$ a multiple of $\begin{bsmallmatrix} 2\\ 1
\end{bsmallmatrix}$? \blank{1.4cm} \; So $\begin{bsmallmatrix} 2\\ 1 \end{bsmallmatrix}$ is
\blank{2.6cm} an eigenvector.
\end{notesbox}

\begin{notesbox}{4. Why It Matters --- and the Road Ahead}
Along an eigenvector, a matrix does nothing but \emph{scale} --- it acts like the single number
$\lambda$. That is what makes these directions so useful (they turn hard matrix problems into easy
number problems). One special case ties back to Unit~5:
\[
  \text{if } A\xx=\zero=0\cdot\xx \text{ for a nonzero }\xx,\ \text{then }\lambda=\blank{0.9cm}
  \text{ and } A \text{ is }\blank{2.4cm}\ (\det A=0).
\]

\vspace{2pt}
\textbf{How do we \emph{find} the $\lambda$'s} without guessing? We solve $\det(A-\lambda I)=0$ ---
a Unit~5 determinant that becomes a quadratic in $\lambda$. That is Lesson~6.1.

\vspace{2pt}
\textbf{The road ahead.} \textbf{6.1} \emph{finding} eigenvalues from $\det(A-\lambda I)=0$\; $\bullet$\;
\textbf{6.2} \emph{diagonalizing} a matrix using its eigenvectors\; $\bullet$\; \textbf{6.3} symmetric
matrices\; $\bullet$\; \textbf{6.4} an application.
\end{notesbox}

\begin{practicebox}
Show your work. To test $\xx$: compute $A\xx$ and check whether it equals $\lambda\xx$ for some number $\lambda$.
\begin{enumerate}[leftmargin=1.6em, itemsep=4pt]
  \item For $A=\begin{bsmallmatrix} 3 & 0\\ 0 & 5 \end{bsmallmatrix}$, compute $A\begin{bsmallmatrix} 1\\ 0 \end{bsmallmatrix}$. Is $\begin{bsmallmatrix} 1\\ 0 \end{bsmallmatrix}$ an eigenvector? What is $\lambda$? \writeline
  \item For $A=\begin{bsmallmatrix} 2 & 1\\ 1 & 2 \end{bsmallmatrix}$ and $\xx=\begin{bsmallmatrix} 1\\ 1 \end{bsmallmatrix}$, find $\lambda$. \writeline
  \item Is $\begin{bsmallmatrix} 1\\ 0 \end{bsmallmatrix}$ an eigenvector of $\begin{bsmallmatrix} 2 & 1\\ 1 & 2 \end{bsmallmatrix}$? Compute $A\begin{bsmallmatrix} 1\\ 0 \end{bsmallmatrix}$ and explain. \writeline
\end{enumerate}
\end{practicebox}

\end{document}
